\section{The completed coheight-one tower}
\label{sec:completed-tower}

We now construct the geometric and spectral comparison objects used by
the detector. The spectral output is a pair of equivariant maps from the
two Morava theories to a completed height-two theory. We also need descent for each
coefficient module and ordinary \(K(2)\)-locality of every profinite
evaluation. These properties allow the later coefficient calculation to
detect the canonical localization map. The inputs
from Barthel--Mann--Ray--Schlank--Senger--Weinstein--Zhou include the
characteristic-\(p\) tower, its cohomology calculations, and the tensor
and completion constructions
in the version dated 26 August 2026 \cite{BMR2026}. We first establish
the connected target and its coefficient descent.
We then construct the map from height three by splitting a completed
connected--\'etale group. The detailed oriented deformation argument is
given in \cref{app:oriented-projection}.

\subsection{The characteristic-\(p\) tower}

Recall that $k=\F_{p^6}$ and that the height-$n$ Honda groups $\Gamma_n$ are over $k$,
for $n=2,3$.  All their geometric endomorphisms are defined over $k$.
Let $D_n$ be the division algebra of invariant $1/n$ over $\Q_p$, with
the convention identifying
\[
 P_n=\Aut_k(\Gamma_n)=\cO_{D_n}^{\times},\qquad
 N_n:P_n\longrightarrow\Z_p^{\times}
\]
with its reduced norm.  Thus $P_n$ is the nonextended stabilizer.  Set
\[
 P_n^1=N_n^{-1}(\mu_{p-1}),\qquad
 \overline Z=P_3/P_3^1\simeq\Z_p,\qquad
 G=P_3\times P_2,\qquad H=P_2\times\overline Z.
\]
The identification of $\overline Z$ with $\Z_p$ will not require a choice
of generator.  We write $\mathcal E_n^0=\cO(1/n)$ for the corresponding
rank-$n$, degree-one Fargues--Fontaine bundle.  A $P_n$-structured form
$\mathcal E_n$ means a form whose torsor of frames has been reduced from
$D_n^{\times}$ to $\cO_{D_n}^{\times}$.  Equivalently, its determinant
framing is specified up to $\Z_p^{\times}$; this is determinant lattice
data, not a lattice in a rank-$n$ rational local system.

The height-exactly-two stratum of the height-three deformation space is
the Laurent field
\[
 K=k((u_2)),\qquad p=u_1=0,\quad u_2\ne0.
\]
Let $L/K$ be the Gross--Torii extension classifying the isomorphisms
$(\Gamma_2)_L\simeq(\Gamma^{\mathrm{un}}_{3,L})^{\circ}$, and let
$\widehat L$ be its valued-field completion.  The two change-of-framing
actions give the continuous $G$-action on $\widehat L$.  Fix compatible
determinant identifications for the two Honda isocrystals, and hence an
identification
$\iota:\det\mathcal E_3^0\simeq\det\mathcal E_2^0$.
These choices exist over $k$: the determinant identifications in
\cite[\S2.7]{BMR2026} require $\F_{p^2}\subset k$.

\begin{proposition}[The completed tower]\label{prop:tower}
The field $\widehat L$ is perfectoid of characteristic $p$, and is
isomorphic, as a valued field over $k$, to the completed perfection of a
Laurent-series field:
\[
 \widehat L\simeq
 \left(\bigcup_{r\ge0}k((z^{1/p^r}))\right)^{\wedge}.
\]
There is a $G$-equivariant identification of v-sheaves over $\Spd k$
\[
 \Spd\widehat L\simeq
 \operatorname{Surj}(\mathcal E_3^0,\mathcal E_2^0).
\]
Consequently
\[
 Y=[\Spd\widehat L/G]
\]
classifies surjections $\mathcal E_3\twoheadrightarrow\mathcal E_2$
between $P_n$-structured forms.  In particular $Y$ carries the universal
surjection used below.

Let $S_3=\ker N_3$ and $T_3=P_3/S_3=\Z_p^{\times}$.  Then
\[
 [\Spd\widehat L/S_3]
 \simeq\operatorname{BC}(-1/2)^*
\]
as v-sheaves, with the following residual $P_2\times T_3$-action.
On extension classes in $\operatorname{Ext}^1(\mathcal E_2^0,\cO)$,
write $\lambda_*$ for pushout by $\lambda$ on $\cO$ and $b^*$ for
pullback by $b$ on $\mathcal E_2^0$.  For the framing action
$f\mapsto bfg^{-1}$ and $a=N_3(g)$, the action is
\[
 (b,a)\cdot\xi
   =\bigl(aN_2(b)^{-1}\bigr)_*(b^{-1})^*\xi.
\]
In particular, the $T_3$-action is pushout by $a$ on the kernel.
Only this restriction of the action on the extension chart is used below.
\end{proposition}

\begin{proof}
The imported result is \cite[Proposition 2.7.1(1)--(5)]{BMR2026}, whose
hypotheses are $h\ge2$, a perfect field containing
$\F_{p^{h(h-1)}}$, and the coheight-one Gross--Torii extension with the
specified determinant identifications.  Here $h=3$ and
$k=\F_{p^6}$ satisfy those hypotheses exactly.

We recall why its formulation gives all surjections and the indicated
structure groups.  The infinite-level two-tower theorem
\cite[Theorem 2.6.3]{BMR2026}, followed by forgetting the étale framing,
identifies $\Spd\widehat L$ with exact sequences
\[
 0\longrightarrow\cO\longrightarrow\mathcal E_3^0
   \longrightarrow\mathcal E_2^0\longrightarrow0
\]
whose determinant identification is $\iota$.  Forgetting that framing is
the $\Z_p^{\times}$-torsor in the proof of
\cite[Proposition 2.7.1(2)]{BMR2026}.  Given a surjection, its kernel is
the line
$\det\mathcal E_3^0\otimes(\det\mathcal E_2^0)^{-1}$; the fixed
$\iota$ trivializes this line and recovers the sequence.  Thus no
additional condition is imposed on the surjection.  Descent along the
two framing torsors gives exactly the asserted description of $Y$, also
in families.

Quotienting by $S_3$ instead allows the source bundle to vary while
retaining its determinant trivialization.  It therefore classifies
extensions
\[
 0\longrightarrow\cO\longrightarrow\mathcal E
   \longrightarrow\mathcal E_2^0\longrightarrow0
\]
with middle term of slope $1/3$.  By
\cite[Proposition 2.7.1(5)]{BMR2026}, these are precisely the extensions
that are nonsplit at every geometric point, which form
$\operatorname{BC}(-1/2)^*$.

To specify the joint action, let $i_f:\cO\to\mathcal E_3^0$ be the
kernel injection normalized by $\iota$.  If $f'=bfg^{-1}$, the
determinant isomorphisms for the two exact sequences give
\[
 g i_f=i_{f'}\,N_3(g)N_2(b)^{-1}.
\]
Indeed, $g$ acts on the source determinant by $N_3(g)$ and $b$ acts on
the target determinant by $N_2(b)$.  The induced diagram of extensions
therefore has kernel map $N_3(g)N_2(b)^{-1}$ and quotient map $b$.
Functoriality of extensions by pushout and pullback proves the displayed
$P_2\times T_3$-action.  This calculation is an identity of bundle maps
on every test object and commutes with base change.  For $b=1$, it is
pushout by $a=N_3(g)$.  On the extension group, $a_*=(aI_2)^*$,
so this is the action by inverse pullback of the central scalar $a^{-1}I_2$.
The underlying quotient and this $T_3$ restriction
are the parts of \cite[Proposition 2.7.1(5)]{BMR2026} used below; the
joint action includes the target determinant factor just computed.
The field assertions are parts (1) and (4) of the same proposition.
\end{proof}

\subsection{Acyclicity of constants}

Choose $C=\C_p$ and a map $\Spd C\to\Spd k$, and write $Y_C$ for
base change over $\Spd k$.  Put
\[
 \widetilde Y=\Spd\widehat L,\qquad
 \widetilde Y_C=\widetilde Y\times_{\Spd k}\Spd C.
\]
The quotient torsors $\widetilde Y\to Y$ and $\widetilde Y_C\to Y_C$
give classifying maps
\[
 Y\longrightarrow B_kG=[\Spd k/\underline G],\qquad
 Y_C\longrightarrow B_CG=[\Spd C/\underline G].
\]
These are relative classifying stacks; the nonextended actions fix
their respective bases.  If $q_n:G\to P_n$ is projection, the classes
denoted below by $e_{P_n}|_Y$ and $e_{P_n}|_{Y_C}$ mean respectively
$u_{\widetilde Y,G,\mathrm b}(q_n^*e_{P_n})$ and
$u_{\widetilde Y_C,G,\mathrm b}(q_n^*e_{P_n})$.
Their base-change compatibility is
\eqref{eq:constant-cochain-naturality} for
$\widetilde Y_C\to\widetilde Y$, already at each finite level.  Introduce
\[
 X=[\Spd\widehat L/P_3^1],\qquad
 X_C=X\times_{\Spd k}\Spd C;
 \quad Y=[X/H],\quad Y_C=[X_C/H].
\]
The notation $\cO^{\flat}$ denotes the characteristic-$p$ structure
sheaf.  Cohomology with these coefficients is derived solid cohomology;
on the group quotients used here, quotient descent computes it as
continuous derived invariants of the parameterized geometric coefficient
complex.  The comparisons
with integral constant coefficients will be made in
\Cref{prop:witt-comparison}.

\begin{proposition}[Acyclicity of the maps of constants]
\label{prop:constant-acyclicity}
The actual maps of constants induce equivalences
\begin{align}
 k&\xrightarrow{\ \sim\ }R\Gamma(P_3^1,\widehat L),
       \label{eq:finite-constant-map}\\
 C^{\flat}&\xrightarrow{\ \sim\ }
       R\Gamma_v(X_C,\cO^{\flat}).
       \label{eq:geometric-constant-map}
\end{align}
These are equivalences of derived solid algebras with residual $H$-action.
The $H$-action on the left sides is trivial.  The maps are natural for
the change of constants $k\to C^{\flat}$ and for the arithmetic actions
preserving the chosen base data.
\end{proposition}

\begin{proof}
Apply the underlying cohomology calculation of
\cite[Proposition~3.6.9]{BMR2026} to the quotient in \Cref{prop:tower},
retaining only its $T_3$-action.  The map of constants and the vanishing
in the other degrees give a fiber sequence of derived solid $k$-modules
with $T_3$-action
\[
 k\longrightarrow R\Gamma(S_3,\widehat L)
   \longrightarrow k_{\chi_T}[-3].
\]
Here $\chi_T$ denotes the actual character on the line in degree three.
This is the restriction of the argument combining quotient descent and cohomology in
\cite[Theorem~3.6.10]{BMR2026} that is needed here.  For $b=1$, the action
formula in \Cref{prop:tower} is pushout by $a\in T_3$ on the kernel.
On the extension group this equals pullback by $aI_2$, or the
action by inverse pullback of $a^{-1}I_2$.  The reduced norm of $aI_2$
in $D_2^\times$ is $a^2$.  Hence the subgroup
\[
 P_3^1/S_3=\mu_{p-1}\subset T_3
\]
acts on the last term by $\alpha\mapsto\alpha^{2}$ or
$\alpha\mapsto\alpha^{-2}$; either convention gives the same vanishing
used below.  The displayed exponent is also checked in
\cite[Lemma 5.2.5]{BMR2026}.  This is a subgroup of the norm quotient;
it must not be replaced by the central $\mu_{p-1}$ in $P_3$.
Since $p\ge5$, this character is nontrivial.  The averaging idempotent
for the finite group $\mu_{p-1}$ is defined in characteristic $p$, so its
derived invariants are exact and annihilate $k_{\chi_T}$.
Taking these invariants proves that the underlying map in
\eqref{eq:finite-constant-map} is an equivalence.  The map itself is
$H$-equivariant: it is induced by constants from the original $G$-action
on $\widehat L$, followed by $P_3^1$-descent, and $G$ fixes $k$.
The forgetful functor from $H$-equivariant objects detects equivalences,
so this gives the required equivalence with its full residual action.
No splitting of the fiber sequence or joint character formula is needed.

For completeness, the finite-field hypothesis in this import is
substantive.  Its proof uses the solid Frobenius comparison
\cite[Lemma 3.6.3 and Corollary 3.6.8]{BMR2026} for
$\bigl(\bigcup_{r\ge0}k((z^{1/p^r}))\bigr)^{\wedge}$
over a \emph{finite} field $k$.
The geometric base change of this point is a punctured perfectoid ball;
the proof solves coefficient Frobenius minus identity on its convergent functions
and then requires surjectivity in solid modules.  For this precise
function ring, \Cref{lem:continuous-as-section} supplies a continuous
section and hence proves surjectivity on every profinite parameter space.
This justifies the solid enhancement directly, without applying a
linear open-mapping theorem to Frobenius minus identity.  The resulting natural
comparison commutes with the $S_3$-action.  Thus
\Cref{prop:tower} verifies the required hypothesis before finite-field
descent is applied.

For the geometric statement, take a quadratic extension $F/\Q_p$ that
is a maximal subfield of $D_2$.  Proposition 3.1.2 of \cite{BMR2026}
identifies
\[
 \bigl(\operatorname{BC}(-1/2)^*\bigr)_C
 \simeq[\mathcal H^1_{F,C}/F],
\]
where $\mathcal H^1_{F,C}$ is the one-dimensional Drinfeld upper
half-plane for $F$ and the additive group $F$ acts by translations.
This identification is equivariant for the central $\Q_p^{\times}$.
It is distinct from the rank-$n$ Drinfeld space
$\Omega_C^{n-1}$ used in the modification argument.
Proposition 3.5.6 of \cite{BMR2026} computes the solid
$\cO^{\flat}$-cohomology of this quotient: it is $C^{\flat}$ in degree
zero, a character line in degree $[F:\Q_p]+1=3$, and zero otherwise.
By the $b=1$ action formula in \Cref{prop:tower}, this norm-quotient
$\mu_{p-1}$ acts by kernel scalars.  Expressing those as central
pullbacks on the target gives reduced-norm character
$\alpha\mapsto\alpha^2$ or its inverse.  Both are nontrivial, so the
top cohomology has no $\mu_{p-1}$-invariants; the choice of action or
dual convention does not affect the conclusion.
The degree-zero identification is the map of constants.  Applying the
same averaging argument proves \eqref{eq:geometric-constant-map}.

Only central equivariance of the Drinfeld description was needed to
prove this vanishing.  The constant map on $X_C$ itself has the full
residual $H$-equivariance, and an equivariant map is an equivalence when
its underlying map is.  This proves the full assertion, without assuming
a $P_2$-equivariant splitting or a larger equivariance of that chart.
Finally, the actions of $P_2$ and $\overline Z$ are over the fixed base,
so they fix $k$ and $C^{\flat}$ pointwise.  All arrows used above are
induced by structure maps and group descent; their stated naturality
follows.
\end{proof}

\subsection{Completed Lubin--Tate theory and coefficient descent}

Write $E_n=E(k,\Gamma_n)$ for ordinary Morava $E$-theory and
$E_n^{\square}=E^{\square}(k^{\delta},\Gamma_n)$ for its solid
enhancement, where $k^{\delta}$ is the discrete solid ring.  For the
valued field $\widehat L$, its solid ring of parameters is
\[
 \widehat L(T)=C_{\cts}(T,\widehat L)
\]
on profinite sets $T$.  Define
\[
 \widehat B^{\square}
   =E^{\square}(\widehat L,(\Gamma_2)_{\widehat L}).
\]
Here $\pi_t$ of a solid spectrum is its solid homotopy module; $(*)$
will denote evaluation in ordinary spectra at the one-point set.
In this section and in \cref{sec:witt-transfer,sec:detection},
$H^s(A,M)$ denotes the point-valued continuous cohomology
$H^s((R\Gamma(A,M))(*))$ for a solid coefficient module $M$.

\begin{proposition}[The completed theory and its coefficients]
\label{prop:completed-theory}
There are continuous $G$-equivariant maps of solid $E_\infty$-rings
\begin{equation}\label{eq:completed-map}
 E_3^{\square}\longrightarrow\widehat B^{\square}
       \longleftarrow E_2^{\square},
\end{equation}
where $G$ acts through $P_3$ on the first source and through $P_2$ on
the second. The right arrow is canonical connected base change.
The target is even periodic, with
\[
 \pi_0\widehat B^{\square}(*)\simeq W(\widehat L)[[u'_1]].
\]
Invariantly, its coefficients are complete base changes of the
height-two deformation ring and its invertible even coefficient lines.
The natural map $W(\widehat L)\to\pi_0\widehat B^{\square}$ is
$G$-equivariant.  Every profinite evaluation
$\widehat B^{\square}(T)$ is an ordinary $K(2)$-local spectrum.

The second map in \eqref{eq:completed-map} induces an $H$-equivariant
equivalence
\[
 E_2^{\square}\xrightarrow{\ \sim\ }
       (\widehat B^{\square})^{hP_3^1},
\]
where $\overline Z$ acts trivially on the source.  Moreover, for every
$t\in\Z$ it induces an equivalence on the individual coefficient
complexes
\[
 \pi_tE_2^{\square}\xrightarrow{\ \sim\ }
       R\Gamma(P_3^1,\pi_t\widehat B^{\square}).
\]
In particular, for $s\ge0$ and all $t$ there are natural isomorphisms
\begin{equation}\label{eq:coefficient-descent}
 H^s(H,\pi_tE_2^{\square})
       \xrightarrow{\ \sim\ }
 H^s(G,\pi_t\widehat B^{\square}).
\end{equation}
These comparisons preserve the coefficient maps from Witt constants.

Finally, put $\Delta=\Gal(k/\F_p)$ and
$G_n=P_n\rtimes\Delta$.  The ordinary stabilizer-descent identification
is
\[
 \bigl((E_n^{\square})^{hG_n}\bigr)(*)\simeq L_{K(n)}S,
 \qquad n=2,3,
\]
compatibly with the canonical maps to $E_n$.
\end{proposition}

The following subsections prove the proposition. Its coefficient
statements concern the connected
base-change map from \(E_2\). We prove them independently of the map
from \(E_3\); that map will be attached to the same endpoint by the
oriented splitting construction below.

\subsection{The connected endpoint and its coefficients}

For the Honda pairs over \(k\), the deformation ideals are
\((p,u_1,u_2)\) at height three and \((p,u'_1)\) at height two.
They are finite regular sequences in the respective deformation rings.
The connected completed pair uses the Noetherian pair
\((k,\Gamma_2)\) as a witness: every
\(C_{\cts}(T,\widehat L)\) is perfect, since inverse Frobenius is
continuous, and flat over \(k\). Thus the endpoint formulas of
\cite[Proposition~4.5.5]{BMR2026} apply through this fixed witness.
The finite Koszul argument below explains the reduction through
hypersheafification and totalization in these instances.

The pointwise coefficient formula and the coefficient-line base change
follow from \cite[Lemma 4.2.12 and Proposition 4.5.5]{BMR2026}, applied to
the constant connected height-two group.  The formal parameter $u'_1$
is a coordinate inherited from a deformation coordinate for $E_2$; no
assertion that $P_2$ fixes this parameter is intended.  Witt vectors and
power series carry their limit structures on profinite parameters.
Functoriality of the perfect-residue-ring deformation construction gives
the asserted map of Witt constants.  The same pointwise description,
together with \cite[Proposition 4.4.2]{BMR2026}, proves $K(2)$-locality
of every evaluation.

The same formulas identify the condensed connected theory
\(E^{\mathrm{cond}}(\widehat L,(\Gamma_2)_{\widehat L})\)
with \(\widehat B^\square\). For a profinite set \(T\), its evaluation is
\[
 E\bigl(C_{\cts}(T,\widehat L),
        (\Gamma_2)_{C_{\cts}(T,\widehat L)}\bigr).
\]
The canonical connected
base-pair map defines
\[
 c^\square:E_2^\square\longrightarrow\widehat B^\square.
\]
It is continuous \(G\)-equivariant by functoriality of the connected
base-pair construction: \(P_3\) acts through the base and \(P_2\)
changes the connected framing. In particular \(G\) acts on its
source through \(P_2\). The family naturality is checked in
\cref{app:oriented-projection}, before constructing the oriented section.
The coefficient and locality statements above apply to this endpoint
without using the full completed intermediate.

We give the coefficient argument separately from the spectral one.
Proposition~4.3.12 of \cite{BMR2026} states the following descent result:
if $(R,\Gamma)$ is a Noetherian Luriean pair and
$S^{-1}\to S^{\bullet}$ is an augmented cosimplicial diagram of flat,
relatively perfect $R$-algebras which is a limit diagram in derived
commutative rings, then both
\[
 E(S^{-1},\Gamma)\longrightarrow
       \operatorname{Tot} E(S^{\bullet},\Gamma)
 \quad\text{and}\quad
 \pi_{2j}E(S^{-1},\Gamma)\longrightarrow
       \operatorname{Tot}\pi_{2j}E(S^{\bullet},\Gamma)
\]
are equivalences; the latter totalization is in $D(\Z)$.
We use this statement for fixed witnesses with finite regular deformation
ideals, and give the finite Koszul justification for their reduction
through totalization.
For any one witness $(R_0,\Gamma_0)$, put
$A_0=\pi_0E(R_0,\Gamma_0)$ and $I_0=\ker(A_0\to R_0)$.
Assume the indicated generators of $I_0$ form a finite regular sequence.
Thus $R_0=A_0/I_0$ has a finite free Koszul resolution and is a perfect
$A_0$-module. Each profinite hypercover diagram for that endpoint is
$A_0$-linear by functoriality from its fixed witness. For the final bar diagram this is
$A_0=W(k)[[u'_1]]$ and $I_0=(p,u'_1)$; its maps are $A_0$-linear
because $P_3^1$ acts trivially on the $E_2$ source.
Derived tensor with $R_0$ therefore commutes with all limits in
$D(A_0)$, including each of these totalizations.

By \cite[Lemma~4.2.12]{BMR2026}, the coefficient objects are derived
$I_0$-complete and their derived reductions are the characteristic-$p$
terms in the augmented diagram. The comparison cone is again derived
complete, since complete objects are closed under limits and cofibers.
Its reduction modulo $I_0$ vanishes by the characteristic-$p$ limit
diagram. If $I_0=(x_1,\ldots,x_c)$, each quotient by
$(x_1^m,\ldots,x_c^m)$ has a finite filtration with subquotients finite
free over $R_0$. The cone therefore vanishes modulo all these quotients.
Their ideals are cofinal with powers of $I_0$, so derived completeness
implies that the cone vanishes. Tensoring with an invertible even
coefficient line also preserves limits. This proves the required
individual-coefficient comparison for every even degree at every
endpoint; the odd coefficient modules vanish. The argument uses the
perfect Koszul quotient for the tensor--totalization interchange.

\paragraph{From coefficient descent to spectral descent.}
These coefficient comparisons also give the spectral comparisons in
ordinary spectra. They give descent for each Eilenberg--Mac Lane
layer, because the forgetful functor from $H\Z$-modules preserves limits;
exactness of totalization then gives descent for every finite Postnikov
interval. For an augmented cosimplicial spectrum
$Z^{-1}\to Z^\bullet$, fix $b$ and $q\leq b$. The canonical fiber sequence
\[
 \tau_{[q-1,b]}Z\longrightarrow\tau_{\leq b}Z
       \longrightarrow\tau_{\leq q-2}Z
\]
has a $(q-2)$-coconnective last term, and totalization preserves this
bound. Hence the first arrow induces an isomorphism on $\pi_q$ both on
the augmentation object and after totalization. The finite-interval
comparison therefore proves the degree-$q$ comparison for
$\tau_{\leq b}Z$. Degrees above $b$ vanish on both sides.
Thus $\tau_{\leq b}Z^{-1}\to\operatorname{Tot}\tau_{\leq b}Z^\bullet$
is an equivalence for every $b$. Taking the limit over $b$, using
Postnikov completeness of ordinary spectra and commutation of limits,
proves the spectral comparison without a bounded-below assumption.

\paragraph{Descent along the completed tower.}
Apply the descent result to \((R,\Gamma)=(k,\Gamma_2)\) and the
\(P_3^1\)-bar diagram of \(\widehat L\), augmented by the constants
\(k\); its map of theories is \(c^\square\). To justify the application
in the solid category, evaluate first on an extremally
disconnected profinite parameter set $T$.  Such evaluation is exact on
condensed modules and preserves limits; moreover, the forgetful functor
from derived commutative algebras creates limits and detects equivalences.
Thus \eqref{eq:finite-constant-map} identifies
\[
 C_{\cts}(T,k)\ \simeq\
 \operatorname{Tot}
 C_{\cts}\bigl((P_3^1)^{\bullet}\times T,\widehat L\bigr)
\]
in derived commutative rings.  All displayed rings are perfect and
flat over $k$, hence relatively perfect. The restricted descent argument
therefore applies to this diagram. The endpoint formulas of
Proposition 4.5.5 and the cotensor comparison of Remark 4.5.6 in
\cite{BMR2026}, with the witness verification above, identify its
spectral and coefficient terms with the evaluations of the required
solid objects and their profinite cotensors. Extremally
disconnected parameter sets detect equivalences; thus this proves both
the spectral and individual-coefficient assertions for \(c^\square\).
Odd coefficient modules vanish on both sides.  This also proves the
specialization of \cite[Theorem 5.2.6]{BMR2026}, whose numerical
hypothesis is exactly $p-1\nmid2$.
Every comparison is induced by the same augmented diagram and its
commuting actions. It therefore retains the residual \(H\)-action and
the maps from Witt constants; \(\overline Z\) acts trivially on the
\(E_2\) source.

\subsection{The map from height three}

The connected theory now has the coefficient descent and locality needed
for detection. To map the height-three sphere into it, we start with the
tensor construction and then pass from the full completed group to its
connected summand. The direction of this last map requires a section of
the deformation functor, rather than the forgetful transformation alone.

The uncompleted tensor maps are the starting point. Examples~4.6.2--4.6.3 of
\cite{BMR2026} use the pointwise constructions of Theorems~4.4.6 and~4.4.8,
followed by Construction~4.5.2, to give
\[
 E_3^{\square}\longrightarrow L_{K(2)}^{\square}E_3^{\square}
   \longrightarrow B^{\square}\longleftarrow E_2^{\square},
\]
where
\[
 B^{\square}
 =L_{K(2)}^{\square}
    (E_3^{\square}\otimes_{SW(k)}^{\square}E_2^{\square})
 \simeq E^{\square}(L,\Gamma^{\mathrm{un}}_{3,L}).
\]
Here \(SW(k)\) is the spherical Witt-vector ring and
\(L_{K(2)}^{\square}\) is localization in solid spectra. These maps have
the projection actions, as in \cite[(5.2.2)--(5.2.3)]{BMR2026}.
The tensor construction uses the constant pair \((k,\Gamma_2)\), whose
base is Noetherian and perfect and whose group is connected of height
two, as required in \cite[Construction~4.4.7]{BMR2026}. The fixed-witness
descent argument above identifies the required hypersheafifications once
the following regular quotient is checked.

For the localized $K$-pair and its uncompleted ind-\'{e}tale $L/K$
extension, the witness has classical deformation ring
\[
 A_K=\bigl(W(k)[[u_1,u_2]][u_2^{-1}]\bigr)^{\wedge}_{(p,u_1)},
 \qquad A_K/(p,u_1)=K.
\]
The sequence $(p,u_1)$ remains regular under localization and this
Noetherian completion. The localized witness is Noetherian Luriean
by \cite[Proposition~4.4.4 and Remark~4.2.2]{BMR2026}; the displayed
$p$-basis below also verifies that $K$ is $F$-finite. The profinite witness hypotheses are
\cite[Propositions~4.5.9, 4.5.10, and 4.5.12, Definition~4.5.11,
and Examples~4.6.1--4.6.3]{BMR2026}.
Concretely, $K(T)=C_{\cts}(T,K)$ is flat over the field $K$ and
relatively perfect: the decomposition
$K=\bigoplus_{j=0}^{p-1}u_2^jK^p$ has continuous coefficient-extraction
maps and hence applies to continuous functions on $T$.
Compactness of $T$ supplies a uniform denominator for inversion of
$u_2$.  This verifies the local-field instance directly; no commutation
of arbitrary localization with totalization is required.
The ind-\'{e}tale $L$-extension retains the same witness as in
Proposition~4.5.12.

\begin{lemma}[The completed oriented projection]
\label{lem:completed-oriented-map}
Put \(F=\widehat L(*)\), regarded here as the ordinary completed field.
For a nonempty profinite set \(T\), put \(R_T=C_{\cts}(T,F)\) and let
\(\Gamma_T=\Gamma^{\mathrm{un}}_{3,R_T}\) be the full completed
\(p\)-divisible group. Let \(C_T\) be the base change of the connected
term in the uniquely split sequence over \(F\). It is identified with
\((\Gamma_2)_{R_T}\) by the fixed tautological Gross--Torii formal
isomorphism. There is a continuous adic \(E_\infty\)-map
\[
 r_T:E(R_T,\Gamma_T)\longrightarrow E(R_T,C_T).
\]
The maps are natural for profinite pullbacks and every continuous
\(G=P_3\times P_2\) family of pair maps. With the terminal value at the
empty parameter set, pointwise application and hypersheafification give
a continuous \(G\)-equivariant map of condensed \(E_\infty\)-rings
\[
 r^{\mathrm{cond}}:
 E^{\mathrm{cond}}(\widehat L,\Gamma^{\mathrm{un}}_{3,\widehat L})
 \longrightarrow E^{\mathrm{cond}}(\widehat L,C_{\widehat L}).
\]

Let \(b_T:E_2\to E(R_T,\Gamma_T)\) be the pointwise ordinary
\(E_2\) tensor-factor map followed by completion, and let
\(c_T:E_2\to E(R_T,C_T)\) be the canonical connected map induced by
\(k\to R_T\). Then \(r_Tb_T\simeq c_T\), naturally for the same
parameters and pair families. This homotopy preserves the connected
residue tag and orientation, and hence its coefficient and even-line maps.
After condensation and solid localization, it gives
\[
 r^{\mathrm{cond}}b^\square\simeq c^\square
\]
as continuous \(G\)-equivariant maps of underlying solid spectra, where
\(b^\square\) is the tensor-factor map followed by completion and
\(c^\square\) is the canonical connected base-change map.
\end{lemma}

\begin{proof}[Construction; full proof in \cref{app:oriented-projection}]
Fix a parameter set \(T\). Write \(\mathscr D_\Gamma^{\mathrm{or}}\)
and \(\mathscr D_C^{\mathrm{or}}\) for oriented deformations with the
specified residue tags of the full and connected reference groups.
The connected--\'{e}tale sequence is uniquely split over the perfect
field \(F\). Let \(Q_T\) be the base change of its \'{e}tale quotient,
so \(\Gamma_T=C_T\oplus Q_T\). An oriented deformation \(C_A\)
over a complete adic test ring \(A\) has a
contractible space of lifts of the tagged \'{e}tale quotient \(Q_T\).
With such a marked lift \(Q_A\), send it to
\(C_A\oplus Q_A\), retaining its connected orientation. This constructs
\(\mathfrak s\) in the left diagram:
\begin{equation}\label{eq:oriented-section-diagram}
\begin{tikzcd}[column sep=large]
 \mathscr D_C^{\mathrm{or}}
    \arrow[r,bend left=18,"\mathfrak s"]
 &\mathscr D_\Gamma^{\mathrm{or}}
    \arrow[l,bend left=18,"\mathfrak f"]
\end{tikzcd}
\qquad
\begin{tikzcd}[column sep=large]
 E(R_T,C_T)\arrow[r,bend left=18,"i_T"]
 &E(R_T,\Gamma_T)\arrow[l,bend left=18,"r_T"] .
\end{tikzcd}
\end{equation}
Here \(\mathfrak f\) takes the connected deformation, and
\(\mathfrak f\mathfrak s\simeq\mathrm{id}\). Corepresentability
reverses the arrows and gives \(r_Ti_T\simeq\mathrm{id}\).
In particular the required map is \(r_T\), supplied by the split
section. The forgetful map alone supplies \(i_T\) in the other direction.

The tensor-factor identity keeps track of which connected map results.
The tensor construction compares the two oriented formal groups through
their common Quillen formal group. On the split deformation its comparison
is the inverse of the connected inclusion. The resulting connected tag
and orientation are therefore exactly those defining \(c_T\), so
\(r_Tb_T\simeq c_T\). The appendix verifies this cancellation on the
actual tags and proves its compatibility with adic refinement and every
continuous group-family map. It then upgrades the identity to the
continuous solid maps in the statement. This identifies the right-hand
tensor composite with the canonical map used in coefficient descent.
\end{proof}

Completion of the pair diagrams followed by the constructed map gives
\[
 B^\square\longrightarrow
 E^{\mathrm{cond}}(\widehat L,\Gamma^{\mathrm{un}}_{3,\widehat L})
 \xrightarrow{r^{\mathrm{cond}}}\widehat B^\square.
\]
The completion map is \cite[Example~4.6.4]{BMR2026}; the full completed
pair uses the perfect-field witness verified in
\cref{app:oriented-projection}. The composite from \(E_3^\square\)
is therefore an actual continuous equivariant ring map. Take this
composite and \(c^\square\) as the arrows in \eqref{eq:completed-map}.
The later detection argument uses this unital equivariant map from
\(E_3^\square\) and the coefficient descent for \(c^\square\).
The identity \(r^{\mathrm{cond}}b^\square\simeq c^\square\) also
identifies the right arrow with the \(E_2\) tensor composite.
This construction uses no coefficient formula
or locality assertion for the nonconnected intermediate, and does not
assert that \(K\to\widehat L\) is relatively perfect.

\subsection{The ordinary sphere source}

It remains to identify the source and the sphere unit after evaluation.
We first spell out the effect of hypersheafification at the point. Let \(a_{\mathrm{hyp}}\) be
hypersheafification on profinite sets with the finite jointly surjective
topology. Every covering sieve of the one-point set is maximal: a
nonempty member of a covering family has a point and hence a section.
Consequently associated-sheaf formation for abelian presheaves leaves
their point value unchanged. The unit
\(P\to a_{\mathrm{hyp}}P\) for a spectrum-valued presheaf is an
isomorphism on associated homotopy sheaves. Evaluating these sheaves at
the point therefore gives an isomorphism on every ordinary homotopy
group, and hence
\[
 P(*)\xrightarrow{\sim}(a_{\mathrm{hyp}}P)(*).
\]
This proof applies to the accessible presheaves on full profinite sets
used here and is natural in the unit. It is the full-profinite version of
\cite[Lemma~13.5 and Remark~13.6]{Clausen2025}; the discrete spectra are
solid by Example~13.4(1) there.

Definition~4.5.7 of \cite{BMR2026} computes condensed localization by
pointwise ordinary localization followed by hypersheafification, and its
restriction to solid spectra is solid localization. Discrete condensation
is the hypersheafification of the constant presheaf, so
\(V^\delta(*)\simeq V\). Thus, naturally in an ordinary spectrum \(V\)
and its maps,
\[
 \bigl(L_{K(n)}^\square(V^\delta)\bigr)(*)\simeq L_{K(n)}V,
\]
with the evaluated localization unit equal to the ordinary unit.
In particular the input \(V=S\) is the ordinary sphere before any
completion. Corollary~4.6.8 of \cite{BMR2026} constructs
\(L_{K(n)}^\square S^\delta\to E_n^\square\) by applying this functor
to the ordinary map \(L_nS\to E_n\). The natural point comparison
identifies it with
\[
 L_{K(n)}L_nS\longrightarrow L_{K(n)}E_n.
\]
Since \(S\to L_nS\) is a \(K(n)\)-equivalence and \(E_n\) is
\(K(n)\)-local, this is precisely the factor of the ordinary sphere
unit \(S\to E_n\) through \(L_{K(n)}S\). Lemma~4.5.8 and the endpoint
formula of Proposition~4.5.5 in that paper identify the target at the
point naturally with \(E_n\): the discrete-adjunction map used there is
adjoint at the point to the identity of \(E_n\).

Proposition~4.6.9 of \cite{BMR2026} identifies the \v Cech augmentation
of that same map with
\(L_{K(n)}^\square S^\delta\xrightarrow{\sim}(E_n^\square)^{hG_n}\).
Its composite to \(E_n^\square\) is the map just identified. Evaluation
at the point preserves limits, so this gives the asserted ordinary
fixed-point equivalence with its canonical map to \(E_n\).
The coefficient-field scope is also exact: Example~4.6.1 there allows an
algebraic field containing the endomorphism field. Here
\(k=\F_{p^6}\) is finite and contains both required fields. The Honda
groups descend to \(\F_p\), and all geometric endomorphisms are
\(k\)-rational. The descent identifications from the \(\F_p\)-models split
the pair-automorphism sequences, giving exactly
\(G_n=P_n\rtimes\Gal(k/\F_p)\). Thus the cited solid descent theorem
applies to this enlarged finite field, and the point comparison above
identifies its source and map for the uncompleted sphere \(S\).

\begin{proof}[Proof of \cref{prop:completed-theory}]
The connected endpoint calculation gives the coefficients and ordinary
locality of every evaluation. The tensor construction and the oriented
section supply the left equivariant ring map; the right map is the
canonical connected base change already constructed.
The constants calculation and finite Koszul
argument prove both fixed-point descent and descent for each coefficient
module, with their residual actions and maps from Witt constants.
Taking derived \(H\)-invariants in the extension
\(P_3^1\to G\to H\) gives \eqref{eq:coefficient-descent};
\(\overline Z\) acts trivially on the \(E_2\) source.
The preceding point-evaluation argument proves the ordinary
stabilizer-descent identification with its sphere unit.
\end{proof}

\begin{remark}
The completion in \eqref{eq:completed-map} is essential to the statements
proved here.  The uncompleted half-sphere equivalence is
\cite[Conjectures 1.1.1 and 5.1.13]{BMR2026};
\cite[Proposition 5.1.17]{BMR2026} proves an implication from that
conjecture to coheight-one splitting.  The completed theorem above
supplies an actual target for detection and does not assert an
equivalence from the uncompleted fixed points to that target.
\end{remark}
